\section{Background and Motivation}
\label{sec:background}

\subsection{Kernel Execution Traces}
\label{sec:bg-traces}

Kernel execution traces record low-level system events during program execution, providing comprehensive visibility into process scheduling, memory allocation, file system operations, network I/O, and inter-process communication~\cite{Noferesti2024Enhancing, Njoku2025Kernel-Level, Giraldeau2016Wait, Gebai2018Survey, Naas2021EZIOTracer, Gellé2021Combining}. Unlike application logs, which capture high-level operations at developer-specified points, kernel traces expose fine-grained system behavior across the entire operating system.

Modern Linux systems support kernel tracing through frameworks such as LTTng~\cite{Desnoyers2006The, Beamonte2015Linux}, eBPF~\cite{eBPF}, and \texttt{ftrace}. We use LTTng, which offers high-throughput tracing with low runtime overhead~\cite{Desnoyers2006The}. Each event records a high-resolution timestamp, event type, CPU core, process metadata (thread and process IDs), and event-specific payload fields (return values, file descriptors, scheduling parameters). These attributes are tightly coupled: event types constrain payload fields, CPU affinity reflects scheduler decisions, and timing patterns vary across workload phases. Realistic trace modeling therefore requires capturing multi-dimensional dependencies across event semantics, temporal dynamics, and resource usage.

Despite their diagnostic value, collecting kernel traces at scale is often impractical due to storage overhead, I/O pressure, and privacy concerns. Engineers frequently rely on limited trace samples collected under controlled conditions, which may fail to capture rare execution paths or long-tail behaviors. This motivates synthetic trace generation as a means to augment limited real data.

\subsection{Diffusion Models for Sequence Generation}
\label{sec:bg-diffusion}

Denoising diffusion probabilistic models (DDPMs)~\cite{ho2020denoisingdiffusionprobabilisticmodels} have emerged as powerful generative models capable of learning complex data distributions. Unlike autoregressive models that generate sequences token-by-token, diffusion models operate by iteratively denoising corrupted signals~\cite{Song2020Denoising}, enabling parallel generation and improved modeling of global structure.

The diffusion process consists of two phases. In the \emph{forward process}~\cite{Arriola2025Block, sutskever2014sequencesequencelearningneural}, Gaussian noise is progressively added to data samples over a fixed number of timesteps. In the \emph{reverse process}, a neural network predicts and removes the injected noise at each step, reconstructing samples from the learned distribution. Transformer-based architectures~\cite{vaswani2023attentionneed} are commonly used as denoisers due to their ability to model long-range dependencies via self-attention—a critical capability for kernel traces, where semantically related events (e.g., system call entry/exit pairs, repeated scheduling patterns) may be separated by thousands of intervening events.

\subsection{Motivation for Constraint-Guided Generation}
\label{sec:bg-motivation}

While neural generative models can capture statistical regularities, statistical plausibility alone is insufficient for kernel traces. Trace-driven system analysis tools assume that traces respect fundamental execution invariants imposed by the operating system, scheduler, and hardware. Violating these invariants yields semantically impossible traces unsuitable for debugging, validation, or training.

Common violation classes include: (i) \emph{invalid event transitions} (e.g., system call exit without entry), (ii) \emph{temporal violations} (inter-event delays outside feasible ranges), and (iii) \emph{attribute inconsistencies} (events on CPU cores or process contexts never observed for those event types). Even infrequent violations can mislead downstream tools or mask true system behaviors.

Relying solely on generative models to implicitly learn such constraints is brittle, particularly with limited or skewed training data. \textsc{SynthTrace} adopts a hybrid approach: explicitly mining execution invariants from real traces and enforcing them through post-hoc validation and repair. This design separates statistical diversity from correctness enforcement, allowing the generator to explore diverse behaviors while ensuring outputs remain consistent with real-system semantics—a prerequisite for safe use in downstream workflows.
